% ==========================================================================
% Paper 2 — Section 8: From Binding to Composition. Overcoming the
% Single-Head Bottleneck (TAN-II Sprint 4.3-A, frozen).
% ==========================================================================
\section{From Binding to Composition: Overcoming the Single-Head
Bottleneck}
\label{sec:composition}

Proposition~\ref{prop:singlehead} states the wall: one head, one symbol.
This section adds the two minimal degrees of freedom that cross it---a
second parallel channel and one two-input algebraic node (operator O5 of
§2.4)---and reports the frozen Sprint 4.3-A audits that separate binding
from composition. The central empirical statement is the second orthogonal
axis of the paper: \emph{binding fidelity $\perp$ combiner error}.

\subsection{The dual-query window}
\label{sec:dualwindow}

The frozen dual-query window is $[A,0,B,Q_1,Q_2]$ (window length $W=5$
kept), with events in the first three slots, queries in the last two, and
the shared window mean $\mu=(K_A+K_B+Q_1+Q_2)/5=2.56$ on the canonical
keys. The query pair $(Q_1,Q_2)=(4.5,5.3)$ is frozen by a
direction-targeting derivation, not by fitting: the pair is required to
rotate the two query directions $\theta_c=\omega S_c$ onto the two key
directions ($\approx44.5^{\circ}$ and $\approx62.8^{\circ}$ versus
$45^{\circ}$ and $63.4^{\circ}$), giving frozen routing margins
$m_A=+0.210631046$ and $m_B=+0.261878523$ (the pre-registration's
hand-computed values $0.2118$/$0.2637$ were four-decimal estimates; the
deviation is recorded as Amendment AM6 with the closed forms adopted and
the originals preserved). Channel $c\in\{1,2\}$ is keyed to slot $3+c$;
queries are control signals, never candidates, and carry value zero.

\subsection{The M0--M1--M2 ladder}
\label{sec:ladder8}

Three models instantiate the negative-control chain of §2.5 on the
canonical scenarios S1 $(5,9)$, S2 $(7,-2)$, S3 $(2,6)$, with operators
$f\in\{+,-\}$. The frozen error table (all values exact):

\begin{itemize}[leftmargin=*,noitemsep,topsep=1pt]
\item \emph{M0 (single winner).} Delivers exactly one operand:
$E_{\mathrm{bind},B}=0$, $E_{\mathrm{bind},A}=|V_B-V_A|>0$, and
$E_{\mathrm{comp}}=|V_A|>0$ for addition (e.g.\ $5$, $7$, $2$ on
S1/S2/S3); the zero-value trap of §7 is re-audited and caught by the
joint criterion.
\item \emph{M1 (parallel binding).} Both channels bind exactly
($E_{\mathrm{bind},A}=E_{\mathrm{bind},B}=0$), and the output is the
\emph{pair} $(\hat V_A,\hat V_B)$: composition is
\textsc{structurally undefined} because no operator connects the two
channels. The frozen non-compositional scalar projection (second-channel
readout, pre-registered as a control) fails with
$E_{\mathrm{comp}}=|V_B-f(V_A,V_B)|>0$: two correctly bound operands are
\emph{not} a composition.
\item \emph{M2 (parallel binding + one two-input combiner).} The node
$y=f(\hat C_1,\hat C_2)$ achieves $E_{\mathrm{bind},A}=E_{\mathrm{bind},B}
=E_{\mathrm{comp}}=0$ exactly on all three scenarios and both operators
($y=14,5,8$ for $+$; $-4,9,-4$ for $-$): the minimal composition node
converts parallel binding into explicit composition.
\end{itemize}

\subsection{The second orthogonal axis and the separation theorem}
\label{sec:separation}

The ensemble audit (keys i.i.d.\ $U(0.3,2.5)$, values i.i.d.\ $U(-3,3)$,
positions and orders randomized, seeds $20260904$--$06$, $N=10^4$ each)
separates the combiner's intrinsic error from the propagated binding
error. Frozen benchmarks (deterministic $2001^2$ quadrature) establish: per-channel
routing rates $r_A=r_B=0.499750$, both-routed probability $P_{\mathrm{both}}=0.220725$ ($22.07\%$),
and unconditional composition errors $E[E_{\mathrm{comp}}^+]=1.116099621$,
$E[E_{\mathrm{comp}}^-]=2.000999500$. The 30,000-history Monte Carlo simulation (seeds $20260904$--$06$,
$N=10^4$ each) yields empirical pooled means of $E[E_{\mathrm{comp}}^+]=1.107671595$
(seed-level: $1.120449 / 1.074606 / 1.127960$) and $E[E_{\mathrm{comp}}^-]=1.975892767$
(seed-level: $1.989288 / 1.959625 / 1.978765$), matching quadrature benchmarks within sampling variance.
The empirical both-routed subset spans $6{,}608 / 30{,}000$ trials ($22.0267\% \approx 22.03\%$;
$2{,}188 / 2{,}258 / 2{,}162$ per seed).

\begin{figure}[t]
\centering
\includegraphics[width=0.70\textwidth]{fig5_composition}
\caption{The minimal composition architecture and audit (Sprint 4.3-A).
(a) The M0--M1--M2 ladder: one winner, two channels, two channels plus one
two-input combiner. (b) The canonical three-error audit across scenarios
S1/S2/S3: only M2 reaches zero composition error. (c) Subtraction
antisymmetry across the four swap counterfactuals. (d) The zero-value trap:
M0's accidental composition zero is caught by the binding error. (e)
Unconditional composition error equals the routing-error benchmark.
(f) Conditioned on correct binding, the combiner error is exactly
$0.0\mathrm{e}0$.}
\label{fig:composition}
\end{figure}

\begin{proposition}[Binding--combiner error separation]
\label{prop:separation}
Given exact operand delivery ($\hat C_c=V_c$ for both channels), the
frozen combiner computes $f$ exactly, so
$E_{\mathrm{comp}}\mid\text{both-routed}=0$. Empirically the maximum over
$3\times10^4$ histories is exactly $0.0\mathrm{e}0$: the combiner adds
zero error, and every unit of unconditional composition error is the
propagation of an upstream routing error
(\emph{routing error $\to$ wrong operand $\to$ composition error}).
\end{proposition}

Proposition~\ref{prop:separation} is the second orthogonal axis:
\emph{binding fidelity $\perp$ combiner error}. Composition capability and
binding fidelity are independent organizational quantities in this family;
the combiner's quality is a property of the algebraic node, the
unconditional error a property of the geometry upstream of it.

\enlargethispage{2\baselineskip}
\subsection{The eight counterfactuals}
\label{sec:cf8}

All eight frozen counterfactuals passed ($\PassFailInconclusive{}$
verdicts, deterministic):

\begin{enumerate}[leftmargin=*,noitemsep,topsep=1pt]
\item \emph{Query swap.} $(Q_1,Q_2)\to(Q_2,Q_1)$: addition invariant,
subtraction exactly antisymmetric ($y\to-y$).
\item \emph{Value swap.} $(K_A,V_B),(K_B,V_A)$: addition invariant
(correct commutative behavior---scored against the channel-level swap
$(9,5)$ as the pairing check, never as a false negative), subtraction
antisymmetric.
\item \emph{Key swap} (keys exchanged, value pairing kept): addition
invariant, subtraction antisymmetric.
\item \emph{Channel-output swap} at the combiner inputs: addition
invariant, subtraction antisymmetric.
\item \emph{AM5 zero-value trap.} $V_A=0$: M0's composition error
vanishes while its binding error does not---caught by the joint audit;
M2 passes genuinely. \textsc{pass}.
\item \emph{Binding-preserving counterfactual.} Value shift
$(5,9)\to(7,11)$: all three errors stay zero and $y=18$; the legacy
single-query controls ($3.0\to A$, $4.0\to B$) reproduce the frozen 4.2-D
routing. \textsc{pass}.
\item \emph{Binding-breaking counterfactual.} Legacy queries $(3.0,4.0)$
in the dual window mis-route channel 2:
$E_{\mathrm{bind},B}=4$ and $E_{\mathrm{comp}}=4$---the causal chain
\emph{routing $\to$ binding $\to$ operand delivery $\to$ composition} is
demonstrated by its failure mode. \textsc{pass}.
\item \emph{Composition-sensitive relational counterfactual.} Across the
three scenarios the output equals $f$ exactly and subtraction
distinguishes the relations ($-4,9,-4$). \textsc{pass}.
\end{enumerate}

Position randomization (six orderings) is bitwise-invariant ($10^{-12}$),
query leakage is excluded ($Q\cap V=\varnothing$; queries never
candidates; query slots carry value zero), and the query-order swap pass
over the whole ensemble preserves addition and antisymmetrizes subtraction
history by history.

\enlargethispage{2.5\baselineskip}
\subsection{The dynamical carrier and spiking readout audit: Phase P1-B}
\label{sec:carrier_audit}

The static algebraic node of \S\ref{sec:ladder8}--\ref{sec:separation} combines
retrieved operand values ($y = C_1 \pm C_2$) once delivered. However, in an online
neural substrate, operands must be continuously maintained and routed through the
internal dynamical state of the neuron, and transmitted across discrete spiking channels
to downstream integrators. Phase~P1-B subjected this continuous dynamical carrier to a pre-registered
confirmatory audit under Protocol~V2.6.

\paragraph{Continuous carrier dynamics and subthreshold state collisions.}
The continuous-time Temporal Attention Neuron carrier evolves as a three-variable ODE
$\mathbf{z}(t) = [u(t), A(t), r(t)]^T$:
\begin{equation}
\tau_u \dot{u}(t) = -u(t) + I(t), \quad
\tau_A \dot{A}(t) = -A(t) + \sigma_A(u(t)), \quad
\tau_r \dot{r}(t) = -r(t) + A(t) \cdot u(t),
\label{eq:p1b_carrier}
\end{equation}
where $I(t)$ is the incoming current pulse sequence, $\sigma_A(\cdot)$ is the
attention activation with threshold $\theta_A$, and $r(t)$ is the surprise-weighted
carrier. At query arrival time $t_Q$, an affine decoder maps the state to delivered operand values:
$\hat{v}_c = g_{\text{dec}} r(t_Q) + b_{\text{dec}}$ ($g_{\text{dec}} = 0.132630817$, $b_{\text{dec}} = -0.026816484$).
Pre-registered fixture G1.5 proved an analytical subthreshold collision theorem: for any interval where $u(t) < \theta_A$,
$\sigma_A(u(t)) \equiv 0$, forcing unmodulated exponential decay $\dot{A} = -A/\tau_A$ and $\dot{r} = -r/\tau_r$.
Consequently, distinct temporal history sequences that differ only during subthreshold periods collapse to analytically identical state vectors $\mathbf{z}(t_Q)$ (residual $< 10^{-10}$), erasing subthreshold events prior to query arrival.

\paragraph{Confirmatory evaluation: floor-limited comparison and unestablished redundancy.}
The confirmatory audit spanned 4,000 unique histories across 3 post-sequence silence
delays ($\Delta t_{\text{silence}} \in \{0, 5, 10\}\,\text{ms}$) and 8 model variants,
totaling 96,000 model-episodes (192,000 role queries). Protocol~V2.6 \S7.3 registered
prerequisite joint delivery accuracy floors ($E_{\text{bind}} \le 0.05$ on both channels)
of $0.80$ for $N=2$ and $0.40$ for $N=3$. The baseline dynamical carrier M1 (evolving from $A(0)=0$) achieved joint delivery accuracy of only $0.0290$ ($58/2{,}000$) for $N=2$ across all delays, and $0.0275$--$0.0280$ ($55$--$56/2{,}000$) for $N=3$. Failing the prerequisite in every cell, all conditions are classified as \texttt{FLOOR\_LIMITED\_COMPARISON}. Two One-Sided Tests (TOST, margin $\pm 0.05$) showed that M1 and the memory-lesioned constant-gate control M2a ($A(t) \equiv A_0 = 0.061886726$) are statistically equivalent within tight $90\%$ confidence intervals ($N=2$: $[-0.001515, 0.001515]$; $N=3$: $[-0.002978, 0.000841]$ at $0\,\text{ms}$, $[-0.002262, 0.001186]$ at $5, 10\,\text{ms}$). However, because both models reside near zero accuracy, this statistical equivalence is an artifact of mutual baseline failure, \emph{not} evidence that temporal gating dynamics are redundant (\texttt{redundancy\_inferred = False}).

\paragraph{Spiking channel limitation.}
To evaluate neural transmission to downstream algebra, carrier states were converted into spike trains via single-unit leaky integrate-and-fire (LIF) decoders under rate (M4\_Count) and first-spike latency (M4\_Latency) codes. Under the registered limitation criterion $E_{\text{comp}} \ge 0.05$, both decoders exhibited severe breakdown across all cells ($E_{\text{comp}} \in [2.21, 2.57]$). Delay-pooled $95\%$ bootstrap confidence intervals were: $N=2$: M4\_Count $[2.393533, 2.502027]$, M4\_Latency $[2.243582, 2.332606]$; $N=3$: M4\_Count $[2.473506, 2.571918]$, M4\_Latency $[2.212398, 2.296450]$. High episode-level abstention rates ($82.1\%$--$83.0\%$; single-query silence $57.15\%$--$58.60\%$) confirm that single-unit spiking decoders cannot reliably transmit continuous-carrier operands, certifying \texttt{DEMONSTRATED\_CHANNEL\_LIMITATION}.

\subsection{Scope of the rung}
\label{sec:scope8}

The established statement is a sufficiency result inside the frozen mechanism family: \emph{parallel binding plus one minimal two-input combiner suffices for identifiable composition, with zero combiner error conditioned on correct binding}. Crucially, this zero composition error is strictly conditional: across the full ensemble of $30{,}000$ histories ($3\times10^4$), the both-routed subset comprises $6{,}608$ histories ($22.03\%$, matching the $22.07\%$ unguided benchmark), with unconditioned pooled errors of $1.108$ ($+$) and $1.976$ ($-$) reflecting upstream routing.

Furthermore, the consensus determination certified under Protocol~V2.6 establishes that M1 fails the baseline delivery prerequisite across all conditions ($2.75\%$--$2.90\%$ joint accuracy), rendering comparisons floor-limited and leaving memory redundancy unestablished, while tested spiking readouts demonstrate channel limitation for single-unit point-process decoders ($E_{\text{comp}} > 2.2$, $82\%$ abstention). These findings establish \emph{model-specific insufficiency} for the continuous ODE carrier and frozen linear readout, not a universal biological impossibility result or a proven requirement for population coding. Internal composition dynamics (Hypothesis~H3), multi-operand composition, and learned operators remain formally deferred.
